\section{Related Work}
\label{sec:related-work}

\subsection{Tabular Foundation Models and In-Context Learning}

TabPFN~\cite{hollmann2023tabpfn} introduced a transformer pretrained once on
millions of synthetic datasets that performs supervised tabular classification
through in-context learning: the training set is presented as a prompt, and a
single forward pass yields predictions for the entire test set, without any
gradient-based fitting on the target data. The follow-up TabPFN~v2
release~\cite{hollmann2025tabpfnv2} scaled the pretraining prior and
architecture, reporting accuracy that matches or exceeds tuned gradient-boosted
ensembles on datasets of up to ten thousand rows while running orders of
magnitude faster. Because a Prior-Data Fitted Network's predictions are shaped
entirely by the class distribution implicit in its pretraining prior together
with the context it is given, its behavior under a training-time class
distribution that differs from what the prior assumes is a distinct concern
from that of a conventionally trained classifier, and has only recently begun
to receive dedicated attention.

DistPFN~\cite{lee2026distpfn} is, to our knowledge, the first method proposed
specifically to correct this failure mode: a test-time posterior adjustment
that rescales TabPFN's predicted class probabilities to counteract label shift
between the training context and the deployment distribution, with a
temperature-scaled variant (DistPFN-T) that adapts the adjustment strength to
the discrepancy between the observed and reference priors. Both variants
require no retraining or architectural change, which is what makes them
practical correction methods for the same class-imbalance setting that
classical resampling and cost-sensitive techniques were designed for on
conventional classifiers. We treat DistPFN and DistPFN-T as first-class
imbalance-correction methods for TabPFN, on equal footing with the classical
corrections described next for gradient-boosted trees, and evaluate the
official formulation directly against the reference implementation before
using it in our study.

\subsection{Correcting Class Imbalance}

The dominant strategies for class imbalance fall into three families:
data-level resampling, algorithm-level cost adjustment, and, for
in-context tabular models, the test-time posterior corrections described
above. SMOTE~\cite{chawla2002smote} synthesizes new minority-class examples by
interpolating between real minority points and their nearest minority
neighbors in feature space, rather than duplicating existing rows, and remains
one of the most widely adopted correction techniques nearly a quarter-century
after its introduction. Random oversampling, by contrast, corrects the
training distribution by duplicating existing minority rows without
synthesizing new feature vectors; both are implemented as standard,
interchangeable components in imbalanced-learn~\cite{lemaitre2017imbalancedlearn},
the toolbox we use for both corrections in this study. Cost-sensitive
learning~\cite{elkan2001costsensitive} instead leaves the training
distribution untouched and reweights the loss so that minority-class errors
are penalized more heavily than majority-class errors in proportion to the
observed class ratio; XGBoost's \texttt{scale\_pos\_weight}
parameter~\cite{chen2016xgboost}, which we use as the weighted correction in
our XGBoost family, is a direct implementation of this idea. All of these
methods share an implicit assumption that is central to this paper: that the
observed class labels used to compute the resampling ratio, the interpolation
neighborhood, or the cost ratio are themselves trustworthy. None of the
methods above were designed with, or have been systematically evaluated
against, the possibility that a fraction of those labels are wrong.

\subsection{Learning Under Label Noise}

A separate literature studies classification when the observed labels
themselves are corrupted, independently of any class imbalance. Early
theoretical work formalized class-conditional label noise, in which the
probability that a true label is flipped depends only on its true class, and
established conditions under which consistent classifiers can still be
recovered~\cite{natarajan2013noisylabels}. Asymmetric label noise, where the
flip probabilities differ across classes or flip in only one direction, was
shown to require different identifiability and denoising strategies than the
symmetric case~\cite{scott2013asymmetricnoise} -- a distinction that maps
directly onto our two directed corruption mechanisms, which are asymmetric by
construction rather than symmetric random flips. Frénay and
Verleysen~\cite{frenay2014labelnoise} survey sources and taxonomies of label
noise and its effects on classifier robustness prior to the deep learning era,
while Song et al.~\cite{song2022noisylabelssurvey} survey the substantially
larger body of work on training deep networks to be robust to noisy labels,
including loss correction, sample selection, and label correction approaches.
Both surveys treat label noise as a phenomenon to be corrected or tolerated
during training; neither surveys, nor the methods they catalog, examine how
label noise interacts with a class-imbalance correction applied
\emph{alongside} it -- which is the specific interaction this paper studies,
not by proposing a new noise-robust training method, but by auditing whether
existing, widely deployed imbalance corrections remain net-beneficial once
that assumption is violated.

\subsection{The Understudied Intersection: Imbalance Correction Under Noisy Labels}

A smaller body of work has begun to treat class imbalance and label noise as
jointly occurring problems rather than independent ones. Liu et
al.~\cite{liu2024imbalancednoisy} observe that standard loss-based sample
selection for noisy-label training is itself biased against minority classes,
since under-learned tail classes and noise-corrupted classes are both
associated with high training loss and are consequently difficult to tell
apart; they propose a class-balanced sample selection procedure with
confidence-based augmentation to correct this bias during training. This line
of work targets the training-time mechanism for handling noisy labels under
imbalance and proposes a new algorithm for that setting. It does not evaluate
whether the standard \emph{imbalance corrections already in widespread use} --
SMOTE, random oversampling, cost-sensitive reweighting, or, for the newer
class of tabular foundation models, test-time posterior adjustment and
class-balanced downsampling -- continue to help, are neutral, or actively harm
performance once a controlled fraction of the labels they operate on are
themselves wrong. To our knowledge, no existing empirical study holds the
total amount of training evidence fixed across imbalance levels while
independently varying a directed, asymmetric label-error rate and measuring
whether the resulting corrections still improve on doing nothing at all --
the gap this paper's controlled, cross-method audit is designed to fill.
